\documentclass[11pt,a4paper]{scrartcl}

\usepackage{../manostilius_lt}

\title{Pilnojo kvadrato išskyrimas}
\author{Andrius Berniukevičius}

\begin{document}

\maketitle

\section{Kas yra pilnojo kvadrato išskyrimas?}
\begin{definition}
\vocab{Pilnasis kvadratas} -- tai reiškinys, kurį galima užrašyti dvinario kvadratu, t. y. forma $(a+b)^2$ arba $(a-b)^2$.
\end{definition}

Pilnojo kvadrato išskyrimas -- tai daugianario pertvarkymas taip, kad jame atsirastų dvinario kvadratas. Šis būdas remiasi greitosios daugybos formulėmis:
\[
(a+b)^2=a^2+2ab+b^2,
\]
\[
(a-b)^2=a^2-2ab+b^2.
\]



\begin{definition}
\vocab{Pilnojo kvadrato išskyrimas} -- tai trinario pertvarkymas į dvinario kvadratą arba į dvinario kvadrato ir skaičiaus sumą ar skirtumą.
\end{definition}

Kad trinarį galėtume užrašyti dvinario kvadratu, pirmasis ir paskutinis jo nariai turi būti kvadratai, o vidurinis narys turi būti lygus dvigubai jų kvadratinių šaknų sandaugai.

\section{Pilnojo kvadrato atpažinimas}

\begin{property}[Pilnojo kvadrato trinaris]
Jei trinaris yra formos
\[
a^2+2ab+b^2,
\]
tai jis lygus $(a+b)^2$. Jei trinaris yra formos
\[
a^2-2ab+b^2,
\]
tai jis lygus $(a-b)^2$.
\end{property}

\setcounter{theorem}{0}

\begin{example}
Išskirkite pilnąjį kvadratą:
\[
x^2+6x+9.
\]
\textbf{Sprendimas}\\
Pirmasis narys yra $x^2$, todėl $a=x$. Paskutinis narys yra $9=3^2$, todėl $b=3$. Patikriname vidurinį narį:
\[
2ab=2\cdot x\cdot3=6x.
\]
Vidurinis narys sutampa, todėl taikome sumos kvadrato formulę:
\[
x^2+6x+9=x^2+2\cdot x\cdot3+3^2=(x+3)^2.
\]
\textbf{Atsakymas:} $(x+3)^2$.
\end{example}

\begin{example}
Išskirkite pilnąjį kvadratą:
\[
4x^2-12x+9.
\]
\textbf{Sprendimas}\\
Pirmasis narys yra $4x^2=(2x)^2$, o paskutinis narys yra $9=3^2$. Todėl tikriname, ar vidurinis narys yra $-2\cdot2x\cdot3$:
\[
-2\cdot2x\cdot3=-12x.
\]
Vidurinis narys sutampa, todėl gauname skirtumo kvadratą:
\[
4x^2-12x+9=(2x)^2-2\cdot2x\cdot3+3^2=(2x-3)^2.
\]
\textbf{Atsakymas:} $(2x-3)^2$.
\end{example}

\begin{example}
Išskirkite pilnąjį kvadratą:
\[
9a^2+24ab+16b^2.
\]
\textbf{Sprendimas}\\
Pirmasis ir paskutinis nariai yra kvadratai:
\[
9a^2=(3a)^2, \qquad 16b^2=(4b)^2.
\]
Patikriname vidurinį narį:
\[
2\cdot3a\cdot4b=24ab.
\]
Taigi tai yra sumos kvadratas:
\[
9a^2+24ab+16b^2=(3a+4b)^2.
\]
\textbf{Atsakymas:} $(3a+4b)^2$.
\end{example}

\begin{example}
Išskirkite pilnąjį kvadratą:
\[
4a^2+12ab+9b^2.
\]
\textbf{Sprendimas}\\
Pirmasis ir paskutinis nariai yra kvadratai:
\[
4a^2=(2a)^2, \qquad 9b^2=(3b)^2.
\]
Patikriname vidurinį narį:
\[
2\cdot2a\cdot3b=12ab.
\]
Taigi tai yra sumos kvadratas:
\[
4a^2+12ab+9b^2=(2a+3b)^2.
\]
\textbf{Atsakymas:} $(2a+3b)^2$.
\end{example}

\begin{example}
Išskirkite pilnąjį kvadratą:
\[
9m^2n^2-6mny+y^2.
\]
\textbf{Sprendimas}\\
Pirmasis ir paskutinis nariai yra kvadratai:
\[
9m^2n^2=(3mn)^2, \qquad y^2=y^2.
\]
Patikriname vidurinį narį:
\[
-2\cdot3mn\cdot y=-6mny.
\]
Taigi tai yra skirtumo kvadratas:
\[
9m^2n^2-6mny+y^2=(3mn-y)^2.
\]
\textbf{Atsakymas:} $(3mn-y)^2$.
\end{example}

\begin{remark}
Ne kiekvienas trinaris yra pilnojo kvadrato trinaris. Pavyzdžiui, reiškinyje $x^2+5x+9$ pirmasis ir paskutinis nariai yra kvadratai, tačiau $2\cdot x\cdot3=6x$, o ne $5x$. Todėl šio trinaro negalime užrašyti $(x+3)^2$ forma.
\end{remark}

\section{Pilnojo kvadrato sudarymas}

Kartais trūksta tik paskutinio nario, kad trinaris taptų pilnojo kvadrato trinariu. Tuomet reikiamą narį pridedame ir iš karto atimame. Reiškinio reikšmė dėl to nesikeičia.

\begin{example}
Išskirkite pilnąjį kvadratą:
\[
x^2+8x.
\]
\textbf{Sprendimas}\\
Pirmasis narys yra $x^2$, o vidurinis narys $8x$ turi būti lygus $2\cdot x\cdot b$. Randame $b$:
\[
2\cdot x\cdot b=8x, \qquad b=4.
\]
Taigi reikia pridėti $4^2=16$. Kad reiškinio reikšmė nepasikeistų, tą patį skaičių ir atimame:
\[
\begin{aligned}
x^2+8x
&=x^2+8x+16-16\\
&=(x+4)^2-16.
\end{aligned}
\]
\textbf{Atsakymas:} $(x+4)^2-16$.
\end{example}

\begin{example}
Išskirkite pilnąjį kvadratą:
\[
x^2-10x+7.
\]
\textbf{Sprendimas}\\
Viduriniam nariui $-10x$ gauti reikia $b=5$, nes $-2\cdot x\cdot5=-10x$. Pridedame ir atimame $5^2=25$:
\[
\begin{aligned}
x^2-10x+7
&=x^2-10x+25-25+7\\
&=(x-5)^2-18.
\end{aligned}
\]
\textbf{Atsakymas:} $(x-5)^2-18$.
\end{example}

\begin{example}
\small
Išskirkite pilnąjį kvadratą:
\[
4x^2+12x+1.
\]
\textbf{Sprendimas}\\
Pirmasis narys yra $(2x)^2$. Vidurinį narį gauname pagal formulę $2\cdot2x\cdot b=12x$, todėl $b=3$. Pridedame ir atimame $3^2=9$:
\[
\begin{aligned}
4x^2+12x+1
&=4x^2+12x+9-9+1\\
&=(2x+3)^2-8.
\end{aligned}
\]
\textbf{Atsakymas:} $(2x+3)^2-8$.
\end{example}

\begin{example}
Išskirkite pilnąjį kvadratą:
\[
x^2-\frac{3}{5}x.
\]
\textbf{Sprendimas}\\
Vidurinio nario koeficiento pusė yra $\frac{3}{10}$, todėl pridedame ir atimame $\left(\frac{3}{10}\right)^2$:
\[
\begin{aligned}
x^2-\frac{3}{5}x
&=x^2-2\cdot\frac{1}{2}\cdot\frac{3}{5}\cdot x \\
&=x^2-2\cdot x\cdot\frac{3}{10} \\
& =\underbrace{x^2-2\cdot x\cdot\frac{3}{10}+\left(\frac{3}{10}\right)^2}_{\left(x-\frac{3}{10}\right)^2}-\left(\frac{3}{10}\right)^2 \\
& =\left(x-\frac{3}{10}\right)^2-\frac{9}{100}.
\end{aligned}
\]
\textbf{Atsakymas:} $\left(x-\frac{3}{10}\right)^2-\frac{9}{100}$.
\end{example}

\begin{example}
Išskirkite pilnąjį kvadratą:
\[
x^2+7x+4.
\]
\textbf{Sprendimas}\\
Vidurinio nario koeficiento pusė yra $\frac{7}{2}$, todėl pridedame ir atimame $\left(\frac{7}{2}\right)^2$:
\[
\begin{aligned}
x^2+7x+4
&=x^2+2\cdot\frac{1}{2}\cdot7x+4 \\
&=x^2+2\cdot x\cdot\frac{7}{2}+4 \\
&=\underbrace{x^2+2\cdot x\cdot\frac{7}{2}+\left(\frac{7}{2}\right)^2}_{\left(x+\frac{7}{2}\right)^2}-\left(\frac{7}{2}\right)^2+4 \\
&=\left(x+\frac{7}{2}\right)^2-\frac{49}{4}+4 \\
&=\left(x+\frac{7}{2}\right)^2-\frac{33}{4}.
\end{aligned}
\]
\textbf{Atsakymas:} $\left(x+\frac{7}{2}\right)^2-\frac{33}{4}$.
\end{example}

\begin{example}
Išskirkite pilnąjį kvadratą:
\[
x^2+3x+1.
\]
\textbf{Sprendimas}\\
Vidurinio nario koeficiento pusė yra $\frac{3}{2}$, todėl pridedame ir atimame $\left(\frac{3}{2}\right)^2$:
\[
\begin{aligned}
x^2+3x+1
&=x^2+2\cdot x\cdot\frac{3}{2}+1\\
&=\underbrace{x^2+2\cdot x\cdot\frac{3}{2}+\left(\frac{3}{2}\right)^2}_{\left(x+\frac{3}{2}\right)^2}-\left(\frac{3}{2}\right)^2+1\\
&=\left(x+\frac{3}{2}\right)^2-\frac{9}{4}+1\\
&=\left(x+\frac{3}{2}\right)^2-\frac{5}{4}.
\end{aligned}
\]
\textbf{Atsakymas:} $\left(x+\frac{3}{2}\right)^2-\frac{5}{4}$.
\end{example}

\begin{example}
Išskirkite pilnąjį kvadratą:
\[
2x^2+\frac{1}{5}x+3.
\]
\textbf{Sprendimas}\\
Iškeliame $2$ prieš pirmus du narius. Skliaustuose vidurinio nario koeficiento pusė yra $\frac{1}{20}$:
\[
\begin{aligned}
2x^2+\frac{1}{5}x+3
&=2\left(x^2+\frac{1}{10}x\right)+3\\
&=2\left(\underbrace{x^2+2\cdot x\cdot\frac{1}{20}+\left(\frac{1}{20}\right)^2}_{\left(x+\frac{1}{20}\right)^2}-\left(\frac{1}{20}\right)^2\right)+3\\
&=2\cdot\left(\left(x+\frac{1}{20}\right)^2-\frac{1}{400}\right)+3\\
&=2\left(x+\frac{1}{20}\right)^2-\frac{1}{200}+3\\
&=2\left(x+\frac{1}{20}\right)^2+\frac{599}{200}.
\end{aligned}
\]
\textbf{Atsakymas:} $2\left(x+\frac{1}{20}\right)^2+\frac{599}{200}$.
\end{example}

\begin{example}
Išskirkite pilnąjį kvadratą:
\[
25a^2b^2+10abc+a.
\]
\textbf{Sprendimas}\\
Pridedame ir atimame $c^2$, kad pirmi trys nariai sudarytų pilnojo kvadrato trinarį:
\[
\begin{aligned}
25a^2b^2+10abc+a
&=25a^2b^2+10abc+c^2-c^2+a\\
&=(5ab+c)^2-c^2+a.
\end{aligned}
\]
\textbf{Atsakymas:} $(5ab+c)^2-c^2+a$.
\end{example}

\section{Uždaviniai}

\begin{problem}
Išskirkite pilnąjį kvadratą:
\begin{tasks}[label={(\arabic*)}, label-offset=0.9em](2)
	\task $x^2+10x+25$
	\task $x^2-14x+49$
	\task $4a^2+20a+25$
	\task $9y^2-24y+16$
	\task $25m^2+30mn+9n^2$
	\task $16p^2-40pq+25q^2$
	\task $a^2-6ab+9b^2$
	\task $x^2+12xy+36y^2$
\end{tasks}
\end{problem}

\begin{problem}
Išskirkite pilnąjį kvadratą:
\begin{tasks}[label={(\arabic*)}, label-offset=0.9em](2)
	\task $9a^2b^2-6abx+x^2$
	\task $4m^2n^2-12mny+9y^2$
	\task $25p^2q^2-20pqr+4r^2$
	\task $x^2-14xy+49y^2$
\end{tasks}
\end{problem}

\begin{problem}
Išskirkite pilnąjį kvadratą:
\begin{tasks}[label={(\arabic*)}, label-offset=0.9em](2)
	\task $x^2+4x$
	\task $x^2-6x$
	\task $x^2+12x+5$
	\task $x^2-8x+3$
	\task $4x^2+20x$
	\task $9x^2-30x+4$
	\task $a^2+14a-2$
	\task $4y^2-4y+6$
\end{tasks}
\end{problem}

\begin{problem}
Išskirkite pilnąjį kvadratą:
\begin{tasks}[label={(\arabic*)}, label-offset=0.9em](2)
	\task $x^2+x+2$
	\task $x^2+5x-1$
	\task $x^2-3x+4$
	\task $x^2-7x+2$
	\task $x^2+\frac{1}{2}x+1$
	\task $x^2-\frac{3}{2}x+5$
	\task $x^2+\frac{7}{3}x-2$
	\task $x^2-\frac{5}{3}x+3$
\end{tasks}
\end{problem}

\begin{problem}
Išskirkite pilnąjį kvadratą, prieš tai iškeldami pirmojo nario koeficientą:
\begin{tasks}[label={(\arabic*)}, label-offset=0.9em](2)
	\task $2x^2+2x+3$
	\task $5x^2-4x+1$
	\task $-3x^2+9x-2$
	\task $4x^2-x+4$
	\task $-x^2+ \frac{1}{4} x+1$
	\task $2x^2-0,1x-1$
	\task $3x^2+ \frac{2}{3} x-1$
	\task $-4x^2-x+2$
\end{tasks}
\end{problem}

\begin{problem}
Išskirkite pilnąjį kvadratą:
\begin{tasks}[label={(\arabic*)}, label-offset=0.9em](2)
	\task $9a^2b^2+12abc+d$
	\task $16x^2y^2-24xyz-9$
	\task $25m^2+20mn-p$
	\task $4r^2s^2+4rst-uv$
\end{tasks}
\end{problem}

\newpage
\answerssection

\begin{problem} \phantom{text}
\begin{tasks}[label={(\arabic*)}, label-offset=0.9em](2)
	\task $(x+5)^2$;
	\task $(x-7)^2$;
	\task $(2a+5)^2$;
	\task $(3y-4)^2$;
	\task $(5m+3n)^2$;
	\task $(4p-5q)^2$;
	\task $(a-3b)^2$;
	\task $(x+6y)^2$.
\end{tasks}
\end{problem}

\begin{problem} \phantom{text}
\begin{tasks}[label={(\arabic*)}, label-offset=0.9em](2)
	\task $(3ab-x)^2$;
	\task $(2mn-3y)^2$;
	\task $(5pq-2r)^2$;
	\task $(x-7y)^2$.
\end{tasks}
\end{problem}

\begin{problem} \phantom{text}
\begin{tasks}[label={(\arabic*)}, label-offset=0.9em](2)
	\task $(x+2)^2-4$;
	\task $(x-3)^2-9$;
	\task $(x+6)^2-31$;
	\task $(x-4)^2-13$;
	\task $(2x+5)^2-25$;
	\task $(3x-5)^2-21$;
	\task $(a+7)^2-51$;
	\task $(2y-1)^2+5$.
\end{tasks}
\end{problem}

\begin{problem} \phantom{text}
\begin{tasks}[label={(\arabic*)}, label-offset=0.9em](2)
	\task $\left(x+\frac{1}{2}\right)^2+\frac{7}{4}$;
	\task $\left(x+\frac{5}{2}\right)^2-\frac{29}{4}$;
	\task $\left(x-\frac{3}{2}\right)^2+\frac{7}{4}$;
	\task $\left(x-\frac{7}{2}\right)^2-\frac{41}{4}$;
	\task $\left(x+\frac{1}{4}\right)^2+\frac{15}{16}$;
	\task $\left(x-\frac{3}{4}\right)^2+\frac{71}{16}$;
	\task $\left(x+\frac{7}{6}\right)^2-\frac{121}{36}$;
	\task $\left(x-\frac{5}{6}\right)^2+\frac{83}{36}$.
\end{tasks}
\end{problem}

\begin{problem} \phantom{text}

\begin{tasks}[label={(\arabic*)}, label-offset=0.9em](2)
	\task $2\left(x+\frac{1}{2}\right)^2+\frac{5}{2}$;
	\task $5\left(x-\frac{2}{5}\right)^2+\frac{1}{5}$;
	\task $-3\left(x-\frac{3}{2}\right)^2+\frac{19}{4}$;
	\task $4\left(x-\frac{1}{8}\right)^2+\frac{63}{16}$;
	\task $-\left(x-\frac{1}{8}\right)^2+\frac{65}{64}$;
	\task $2\left(x-\frac{1}{40}\right)^2-\frac{801}{800}$;
	\task $3\left(x+\frac{1}{9}\right)^2-\frac{28}{27}$;
	\task $-4\left(x+\frac{1}{8}\right)^2+\frac{33}{16}$.
\end{tasks}
\end{problem}

\begin{problem} \phantom{text}
\begin{tasks}[label={(\arabic*)}, label-offset=0.9em](2)
	\task $(3ab+2c)^2-4c^2+d$;
	\task $(4xy-3z)^2-9z^2-9$;
	\task $(5m+2n)^2-4n^2-p$;
	\task $(2rs+t)^2-t^2-uv$.
\end{tasks}
\end{problem}

\end{document}
